% Introduction
\section{Introduction}
\label{sec:introduction}

Eternal inflation \cite{guth_eternal} produces an ensemble of causally disconnected
or weakly related inflationary domains---a multiverse---for which standard
quantum-field-theory calculations do not prescribe how to compute probabilities.
The central difficulty, known as the measure problem
\cite{freivogel_measure}, is that any enumeration over an infinite population
requires a regularization scheme (a ``measure''), and different schemes yield
different answers.

The simplest regularization---a global census that counts each trajectory
equally---is trajectory-cutoff dependent and generically favors the most
numerous basin rather than any physically motivated reference point.
This is not a prediction; it is an artifact of the counting rule.

We propose \isr (\ISR), a configuration-space locality-weighted unobserved-sector
framework that conditions the ensemble on an observable patch $x_0$ and
weights sectors by their distance to $x_0$ in field space and basin structure.
The weighting uses a kernel $K_\ell[d(x, x_0)]$ that decays with increasing
sector distance, controlled by a scale parameter $\ell$.
ISR is a measure-construction proposal, not a local hidden-variable model.

\subsection{Scope of the proposal}
\label{sec:scope}

It is important to state what ISR does not claim:

\begin{itemize}
  \item ISR does \textbf{not} solve the full eternal-inflation measure problem.
        Volume divergences in an infinite multiverse are not regulated here.
  \item ISR does \textbf{not} regulate spacetime volume divergences.
        It operates on a finite ensemble of stochastically generated trajectories.
  \item The cutoff studied here is a \textbf{trajectory-sampling cutoff},
        not a cosmological volume cutoff.
        It tests whether probability estimates stabilize as more samples are
        drawn from a given stochastic process.
  \item The locality kernel is a \textbf{prior over configuration-space relevance},
        not a claim about causal contact between bubbles.
\end{itemize}

ISR is a conditional weighting scheme inside a generated sector ensemble.
In future work it could be composed with an existing cosmological measure
rather than replacing one.

We implement a stochastic toy landscape with multiple vacuum basins and
simulate 10,000 trajectories, comparing ISR against a naive global census.
We also run 100 sigma-by-seed stress tests to assess robustness, plus
landscape-perturbation tests to verify that results are not an artifact of
the specific well placement.

\begin{quote}
\textbf{Central claim.}
\ISR does not count all inflationary domains as equally relevant.
It conditions the ensemble on an observable patch and weights unobserved sectors
by locality in field space and basin structure.
The proposal succeeds in a toy model only if the resulting measure is more
stable under finite-sample cutoff expansion than a naive global census and
remains robust across reasonable kernels, seeds, and initial-condition spreads.
\end{quote}

\ISR should be judged by stability, robustness, and transparency about
failure cases, not by whether it can be tuned to prefer a desired vacuum.
The toy model demonstrates a computational behavior, not a direct physical
derivation from quantum gravity.
